% Bagian: second-order-logic
% Bab: syntax-and-semantics
% Subbagian: expressive-power

\documentclass[../../../include/open-logic-section]{subfiles}

\begin{document}

\olfileid[id]{sol}{syn}{exp}
\olsection{Daya Ekspresif}

\begin{explain}
Kuantifikasi atas variabel orde kedua menghasilkan peningkatan yang sangat
besar dalam daya ekspresif bahasa dibandingkan dengan logika orde pertama.
Kuantifikasi eksistensial orde kedua memungkinkan kita menyatakan bahwa
terdapat fungsi atau relasi yang mempunyai sifat-sifat tertentu. Dalam logika
orde pertama, satu-satunya cara untuk melakukan hal itu ialah menentukan
suatu simbol nonlogis (yakni, !!a{function} atau !!{predicate}) bagi tujuan
tersebut. Kuantifikasi universal orde kedua memungkinkan kita menyatakan
bahwa semua himpunan bagian dan relasi pada !!{domain}, serta semua fungsi
dari !!{domain} itu ke dirinya sendiri, mempunyai suatu sifat. Dalam
logika orde pertama, kita hanya dapat menyatakan bahwa himpunan bagian,
relasi, atau fungsi yang ditetapkan kepada salah satu simbol nonlogis bahasa
mempunyai suatu sifat. Dan ketika kita menyatakan bahwa terdapat himpunan
bagian, relasi, atau fungsi yang mempunyai suatu sifat, atau bahwa semuanya
mempunyai sifat itu, kita juga dapat menggunakan kuantifikasi orde kedua
dalam menentukan sifat tersebut. Hal ini memungkinkan kita mendefinisikan
relasi yang tidak dapat didefinisikan dalam logika orde pertama dan
mengungkapkan sifat-sifat domain yang tidak dapat diungkapkan dalam logika
orde pertama.
\end{explain}

\begin{defn}
Jika $\Struct{M}$ merupakan !!a{structure} bagi suatu bahasa~$\Lang{L}$,
suatu relasi~$R \subseteq \Domain{M}^2$ \emph{dapat didefinisikan}
dalam~$\Lang{L}$ jika terdapat !!{formula}~$!A_R(x, y)$ yang hanya mempunyai
variabel $x$ dan~$y$ sebagai variabel bebas, sedemikian sehingga $R(a, b)$
berlaku (yakni, $\tuple{a,b} \in R$) jika dan hanya jika
$\Sat{M}{!A_R(x, y)}[s]$ bagi $s(x) = a$ dan $s(y) = b$.
\end{defn}

\begin{ex}
Dalam logika orde pertama kita dapat mendefinisikan relasi
identitas~$\Id{\Domain{M}}$ (yakni, $\Setabs{\tuple{a,a}}{a \in
  \Domain{M}}$) dengan formula $\eq[x][y]$. Dalam logika orde kedua, kita
dapat mendefinisikan relasi ini \emph{tanpa~$\eq$}. Sebab, jika $a$ dan $b$
merupakan !!{element} yang sama dari~$\Domain{M}$, keduanya menjadi
!!{element}s dari himpunan-himpunan bagian yang sama dari~$\Domain{M}$
(karena himpunan ditentukan oleh !!{element}s di dalamnya). Sebaliknya, jika $a$ dan
$b$ berbeda, keduanya bukan !!{element}s dari himpunan-himpunan bagian yang
sama: misalnya, $a \in \{a\}$ tetapi $b \notin \{a\}$ jika $a \neq b$.
Jadi, ``menjadi !!{element}s dari himpunan-himpunan bagian yang sama
dari~$\Domain{M}$'' merupakan relasi yang berlaku atas $a$ dan $b$ jika dan
hanya jika $a = b$. Relasi itu dapat diungkapkan dalam logika orde kedua,
karena kita dapat menguantifikasi atas semua himpunan bagian dari~$\Domain{M}$.
Oleh karena itu, !!{formula} berikut mendefinisikan $\Id{\Domain{M}}$:
\[
\lforall[X][(X(x) \liff X(y))]
\]
\end{ex}

\begin{prob}
Tunjukkan bahwa $\lforall[X][(X(x) \lif X(y))]$ (perhatikan:
$\lif$, bukan~$\liff$!) mendefinisikan $\Id{\Domain{M}}$.
\end{prob}

\begin{ex}
Jika $R$ merupakan !!{predicate} beraritas dua, $\Assign{R}{M}$ merupakan
relasi beraritas dua pada~$\Domain{M}$. Barangkali agak membingungkan, kita
akan menggunakan $R$ baik sebagai !!{predicate}~$R$ maupun sebagai relasi
$\Assign{R}{M}$ itu sendiri. \emph{Penutupan transitif}~$R^+$ dari~$R$ ialah
relasi yang berlaku antara $a$ dan $b$ jika dan hanya jika untuk suatu $c_1$,
\dots, $c_k$, berlaku $R(a,c_1)$, $R(c_1, c_2)$, \dots, $R(c_k,b)$. Ini
mencakup kasus ketika $k = 0$, yakni, jika $R(a,b)$ berlaku, maka $R^+(a,b)$
juga berlaku. Ini berarti $R \subseteq R^+$. Bahkan, $R^+$ merupakan relasi
terkecil yang memuat~$R$ dan bersifat transitif. Dalam logika orde kedua kita
dapat menyatakan bahwa $X$ merupakan relasi transitif yang memuat~$R$:
\begin{multline*}
  !B_R(X) \ident \lforall[x][\lforall[y][(R(x,y) \lif X(x, y))]] \land {}\\
\lforall[x][\lforall[y][\lforall[z][((X(x,y) \land X(y,z)) \lif X(x,
      z))]]].
\end{multline*}
Konjung pertama menyatakan bahwa $R \subseteq X$, dan konjung kedua
menyatakan bahwa $X$ transitif.

Menyatakan bahwa $X$ merupakan relasi terkecil semacam itu berarti menyatakan
bahwa relasi itu sendiri termuat dalam setiap relasi yang memuat~$R$ dan
bersifat transitif. Jadi, kita dapat mendefinisikan penutupan transitif
dari~$R$ dengan !!{formula}
\[
R^+(X) \ident !B_R(X) \land \lforall[Y][(!B_R(Y) \lif
  \lforall[x][\lforall[y][(X(x, y) \lif Y(x,y))]])].
\]
Kita mempunyai $\Sat{M}{R^+(X)}[s]$ jika dan hanya jika $s(X) = R^+$.
Penutupan transitif dari~$R$ tidak dapat diungkapkan dalam logika orde
pertama.
\end{ex}

\end{document}
